\documentclass[11pt,a4paper]{article}
\usepackage[T1]{fontenc}
\usepackage[utf8]{inputenc}
\usepackage{lmodern,amsmath,amssymb,amsfonts,amsthm,mathtools}
\usepackage{graphicx,booktabs,microtype,enumitem,hyperref,xcolor,fancyhdr}
\usepackage[margin=23mm,headheight=16pt]{geometry}
\definecolor{accent}{HTML}{8D443B}
\hypersetup{colorlinks=true,linkcolor=accent,urlcolor=accent,
pdftitle={Chapter 6: Symmetries and Conservation Laws -- 37 Worked Problems}}
\pagestyle{fancy}\fancyhf{}\lhead{\small QM / Symmetries \& Conservation Laws}
\rhead{\small K. Azad \ \textbullet\ Chapter 6}\cfoot{\thepage}
\setlength{\parindent}{0pt}\setlength{\parskip}{5pt}
\allowdisplaybreaks[2]
\numberwithin{equation}{section}
\newcommand{\hb}{\hbar}
\newcommand{\ii}{\mathrm i}
\newcommand{\ee}{\mathrm e}
\newcommand{\dd}{\,\mathrm d}
\newcommand{\ket}[1]{\lvert #1\rangle}
\newcommand{\bra}[1]{\langle #1\rvert}
\newcommand{\braket}[2]{\langle #1\mid #2\rangle}
\newcommand{\av}[1]{\langle #1\rangle}
\newcommand{\T}{\hat T}
\newcommand{\PP}{\hat{\mathcal P}}
\newcommand{\RR}{\hat R}
\newcommand{\UU}{\hat U}
\newcommand{\HH}{\hat H}
\newcommand{\LL}{\hat{\mathbf L}}
\newcommand{\SS}{\hat{\mathbf S}}
\newcommand{\mm}{\hat{\mathbf M}}
\newcommand{\x}{\hat x}
\newcommand{\p}{\hat p}
\newcommand{\problem}[2]{\subsection*{Problem #1: #2}
\addcontentsline{toc}{subsection}{Problem #1: #2}}
\newcommand{\answer}[1]{\par\noindent{\color{accent}\bfseries Result.} #1}
\newcommand{\phys}[1]{\par\noindent{\itshape Physical interpretation.} #1}
\newcommand{\caveat}[1]{\par\noindent{\color{accent}\bfseries Assumption/check.} #1}
\newcommand{\vc}[1]{\hat{\boldsymbol{#1}}}
\title{{\color{accent}\small MATHEMATICAL PHYSICS · INDEPENDENT EDUCATIONAL COMPANION}\\[12pt]
\Huge\bfseries Symmetries \&\\ Conservation Laws\\[10pt]
\Large Quantum Mechanics — Chapter 6\\[9pt]
\normalsize Problems 6.1--6.37, with original mathematical derivations}
\author{Kumar Azad\\
\small \url{https://krazad0.github.io/QM-CH1-SOLUTION/chapter6/}}
\date{November 2026}
\begin{document}
\maketitle
\begin{abstract}
This independent study companion develops the \emph{37 numbered
exercises} of Chapter 6 of Griffiths and Schroeter,
\emph{Introduction to Quantum Mechanics}, third edition
(Cambridge University Press, 2018). Topics include translations
and Bloch theory; conservation and parity; SO(3), SU(2) and
Wigner--Eckart relations; selection rules; propagators;
accidental degeneracy and the quantum Laplace--Runge--Lenz
vector; Galilean boosts; and time reversal/Kramers degeneracy.
Exercise descriptions are independently paraphrased. This is
not an official solutions manual, arXiv preprint, or peer-reviewed
research. Some long exercises are presented as compact
reproducible derivations rather than exhaustive proofs
of all general identities.
\end{abstract}
\tableofcontents
\clearpage
\section{Conventions and symmetry roadmap}
The active one-dimensional translation convention used by
the textbook is
\[
\T(a)f(x)=f(x-a),\qquad
\T(a)=\ee^{-\ii a\p/\hb},\qquad
\T(a)^\dagger \x\T(a)=\x+a.
\]
Spatial inversion and rotations obey
\[
\PP f(\mathbf r)=f(-\mathbf r),\quad
\PP^2=I,\qquad
\RR_{\mathbf n}(\varphi)=
\exp(-\ii\varphi\,\mathbf n\cdot\LL/\hb).
\]
For a time-independent Hamiltonian
$\UU(t)=\exp(-\ii\HH t/\hb)$.
An operator $Q$ with no explicit time
dependence is conserved if $[\HH,Q]=0$.
A parity-even operator connects
states of the same parity; a parity-odd
operator can connect only opposite parities.
Tensor-operator rotational selection rules
are encoded by the Wigner--Eckart theorem.
A spin-half time-reversal transformation
is \emph{antiunitary}, not unitary.
\section{Translations, Bloch waves, momentum: Problems 6.1--6.7}
\problem{6.1}{Parity, mirror reflections, and hydrogenic harmonics}
\textbf{(a)} A reflection in the $xy$ plane sends
$(x,y,z)\mapsto(x,y,-z)$. Following it with a
half turn about $z$ sends
$(x,y,-z)\mapsto(-x,-y,-z)=-\mathbf r$.
Thus spatial inversion can be realized by a
mirror reflection followed by a rotation.
\textbf{(b)} In spherical coordinates
$x=r\sin\theta\cos\phi$, $y=r\sin\theta\sin\phi$,
$z=r\cos\theta$. The transformation
\[
\theta\longmapsto\pi-\theta,\qquad
\phi\longmapsto\phi+\pi
\]
changes each Cartesian coordinate's sign,
and $r$ stays fixed. Consequently
\begin{equation}
\boxed{(\PP\psi)(r,\theta,\phi)=
\psi(r,\pi-\theta,\phi+\pi).}
\end{equation}
\textbf{(c)} For $\psi_{n\ell m}=R_{n\ell}(r)
Y_\ell^m(\theta,\phi)$, use the standard
harmonic $Y_\ell^m=N_{\ell m}
P_\ell^m(\cos\theta)\ee^{\ii m\phi}$.
Since
$P_\ell^m(-u)=(-1)^{\ell+m}P_\ell^m(u)$ and
$\ee^{\ii m\pi}=(-1)^m$,
\[
Y_\ell^m(\pi-\theta,\phi+\pi)
=(-1)^{\ell+2m}Y_\ell^m(\theta,\phi)
=(-1)^\ell Y_\ell^m.
\]
The radial function depends only on $r$:
\begin{equation}
\boxed{\PP\psi_{n\ell m}=(-1)^\ell\psi_{n\ell m}.}
\end{equation}
This conclusion is valid for \emph{any} central
potential, not just the Coulomb potential.
\phys{Even $\ell$ gives even parity; odd $\ell$
gives odd parity.}

\problem{6.2}{Exponential of a Hermitian generator is unitary}
Let $Q^\dagger=Q$ and
$U=\exp(\ii Q)=\sum_{n\ge0}(\ii Q)^n/n!$.
Taking the adjoint termwise,
\[
U^\dagger=\sum_{n\ge0}
\frac{(-\ii Q^\dagger)^n}{n!}
=\exp(-\ii Q).
\]
The two exponentials commute because
they contain the same operator $Q$.
Therefore
\begin{equation}
\boxed{U^\dagger U=
\ee^{-\ii Q}\ee^{\ii Q}=
\ee^{0}=I,\qquad UU^\dagger=I.}
\end{equation}
By the spectral theorem this remains
valid for a self-adjoint (possibly
unbounded) $Q$ through the functional
calculus; a formal power series need
not converge on every state.

\problem{6.3}{Momentum does not change under spatial translations}
Take $\T(a)f(x)=f(x-a)$,
$\p=-\ii\hb\partial_x$, and
$\T(a)^\dagger=\T(-a)$. Acting on
a differentiable test function,
\begin{align}
(\T^\dagger\p\T f)(x)
&=\T(-a)[-\ii\hb f'(x-a)]\\
&=-\ii\hb f'(x)=(\p f)(x).
\end{align}
Hence
\begin{equation}
\boxed{\T(a)^\dagger\p\T(a)=\p.}
\end{equation}
Equivalently $[\p,\T(a)]=0$
because $\T(a)=\exp(-\ii a\p/\hb)$.
\phys{A change of origin shifts position
but not the momentum.}

\problem{6.4}{Translation of an arbitrary operator series}
First, the example discussed in the text gives
$\T^\dagger\x\T=\x+a$; Problem 6.3 gives
$\T^\dagger\p\T=\p$.
For any operators $A,B$ and unitary $T$,
\[
T^\dagger(AB)T=
(T^\dagger AT)(T^\dagger BT)
\]
by inserting $TT^\dagger=I$.
Apply the identity repeatedly:
\[
\T^\dagger \x^m\p^n\T
=(\x+a)^m\p^n.
\]
Therefore if $Q=\sum_{m,n}c_{mn}\x^m\p^n$
in the ordering specified in the exercise,
\begin{equation}
\boxed{\T^\dagger Q(\x,\p)\T
=Q(\x+a,\p).}
\end{equation}
The result extends to ordered products
and suitable functions by functional
calculus, but operator ordering should
be preserved when $[\x,\p]\ne0$.

\problem{6.5}{Proving the periodic factor in Bloch's theorem}
Let the Hamiltonian commute with translation
by $a$, and let its chosen eigenstate obey
\[
\T(a)\psi(x)=\psi(x-a)
=\ee^{-\ii qa}\psi(x).
\]
Define $u(x)=\ee^{-\ii qx}\psi(x)$,
so $\psi(x)=\ee^{\ii qx}u(x)$.
Then
\begin{align}
u(x-a)
&=\ee^{-\ii q(x-a)}\psi(x-a)\\
&=\ee^{-\ii qx}\ee^{\ii qa}
\ee^{-\ii qa}\psi(x)
=u(x).
\end{align}
Shifting $x\mapsto x+a$ shows
$u(x+a)=u(x)$:
\begin{equation}
\boxed{\psi_q(x)=\ee^{\ii qx}u_q(x),
\qquad u_q(x+a)=u_q(x).}
\end{equation}
Here $q$ and $q+2\pi/a$ describe
the same Bloch translation eigenvalue.
A periodic potential conserves \emph{crystal}
momentum, not necessarily ordinary momentum.

\problem{6.6}{The lowest two bands in the periodic finite well}
\textbf{(a) Derivation.}
Write $\psi_q(x)=\ee^{\ii qx}u_q(x)$, where
$u_q(x+a)=u_q(x)$. Differentiating twice,
\[
\psi_q''=\ee^{\ii qx}
(u_q''+2\ii q u_q'-q^2u_q).
\]
Substitute in $-\hb^2\psi_q''/(2m)+V\psi_q=E\psi_q$:
\begin{equation}
\boxed{\left[-\frac{\hb^2}{2m}\frac{\dd^2}{\dd x^2}
-\frac{\ii\hb^2q}{m}\frac{\dd}{\dd x}
+\frac{\hb^2q^2}{2m}+V(x)\right]u_q
=E_q u_q.}
\end{equation}
The opposite Bloch-phase convention changes
the sign of the linear-derivative term but
not the energies.
\textbf{(b) Exact matching reference.}
Use $E_*=\hb^2/(2ma^2)$,
$s=x/a$, $\epsilon=E/E_*$, and
$v=-20$ on $1/4<s<3/4$ (zero outside).
Translation of the origin gives two
segments of length $a/2$ with potentials
$-20E_*$ and $0$. At a constant-potential
segment of dimensionless length $d$,
a state $(\psi,\psi')^T$ propagates through
\begin{equation}
M(k,d)=
\begin{pmatrix}
\cos(kd)&\sin(kd)/k\\
-k\sin(kd)&\cos(kd)
\end{pmatrix},
\qquad k=\sqrt{\epsilon-v}.
\end{equation}
For imaginary $k$ the entries are real
hyperbolic functions. Bloch boundary
conditions imply that the eigenvalues of
$M(\sqrt\epsilon,\tfrac12)
M(\sqrt{\epsilon+20},\tfrac12)$
are $\ee^{\pm\ii qa}$. Since $\det M=1$,
\begin{equation}
\boxed{\cos(qa)=\frac12\operatorname{tr}
\left[M(\sqrt\epsilon,\tfrac12)
M(\sqrt{\epsilon+20},\tfrac12)\right].}
\end{equation}
A root search gives the \emph{two lowest}
energy levels in units of $E_*$:
\begin{equation}
\begin{array}{c|rr}
qa& E_1/E_*&E_2/E_*\\\hline
-\pi&-6.981932&5.618248\\
-\pi/2&-10.060437&13.552029\\
0&-11.989299&28.848077\\
\pi/2&-10.060437&13.552029\\
\pi&-6.981932&5.618248
\end{array}
\end{equation}
\textbf{(c) Reproducible finite differences.}
On an $N$-point periodic grid with
$h=1/N$, approximate the Bloch-covariant
kinetic operator by
\begin{equation}
(H_qu)_j=\frac{
2u_j-\ee^{\ii qh}u_{j+1}
-\ee^{-\ii qh}u_{j-1}}{h^2}
+v(s_j)u_j,
\quad j\bmod N.
\end{equation}
The resulting matrix is Hermitian.
For $N=400$ or $800$, diagonalize
its two lowest eigenvalues for
$qa\in[-\pi,\pi]$; increasing $N$
checks convergence to the exact-matching
table. The plotting script writes
\texttt{figures/bloch\_bands.pdf}.
\phys{The two bands are even in $q$ because
the potential and Hamiltonian are real;
their band gap is an effect of periodic
scattering, not an imposed boundary wall.}

\problem{6.7}{Total momentum conservation for an interacting pair}
Let
$H=\p_1^2/(2m_1)+\p_2^2/(2m_2)+
V(|x_1-x_2|)$.
\textbf{(a)} A simultaneous translation
of both coordinates acts as
\begin{equation}
\T(a)\psi(x_1,x_2)
=\psi(x_1-a,x_2-a)
=\exp[-a(\partial_{x_1}+\partial_{x_2})]\psi.
\end{equation}
Using $p_j=-\ii\hb\partial_{x_j}$,
\[
\boxed{\T(a)=\exp(-\ii aP/\hb),\qquad
P=p_1+p_2.}
\]
\textbf{(b)} The kinetic energy
commutes with $P$.
With $r=x_1-x_2$,
\[
[P,V(|r|)]=-\ii\hb
(\partial_{x_1}+\partial_{x_2})V(|r|)
=0
\]
since $\partial_{x_1}V=V'(r)$ and
$\partial_{x_2}V=-V'(r)$ wherever
the derivative is defined. Therefore
\begin{equation}
\boxed{[H,P]=0,\qquad
\frac{\dd\av P}{\dd t}=0.}
\end{equation}
\phys{Internal forces can exchange
momentum between the particles but
cannot change the total momentum.}

\section{Parity and rotation: Problems 6.8--6.17}
\problem{6.8}{Hermiticity and spectrum of parity}
\textbf{(a)} In one dimension
$(\PP g)(x)=g(-x)$, and
\[
\braket f{\PP g}=\int_{-\infty}^\infty
f^*(x)g(-x)\dd x.
\]
Set $u=-x$, reverse integration limits:
\[
\braket f{\PP g}=
\int_{-\infty}^\infty f^*(-u)g(u)\dd u
=\braket{\PP f}g.
\]
Therefore $\boxed{\PP^\dagger=\PP}$.
The same coordinate substitution works in
three dimensions because the Jacobian's
absolute value equals one.
\textbf{(b)} Inverting coordinates twice gives
$\PP^2=I$. If $\PP\psi=\lambda\psi$,
then $\PP^2\psi=\lambda^2\psi=\psi$.
For a nonzero eigenstate
\begin{equation}
\boxed{\lambda^2=1,\qquad\lambda=\pm1.}
\end{equation}
Moreover $\PP^\dagger\PP=\PP^2=I$,
so parity is both Hermitian and unitary.

\problem{6.9}{Commutator and anticommutator parity classification}
For an operator $A$, let
$A'=\PP^\dagger A\PP=\PP A\PP$.
If $A$ is \emph{even}, $A'=A$.
Multiply on the left by $\PP$
and use $\PP^2=I$: $A\PP=\PP A$.
Hence
\[
\boxed{[\PP,A]=0\qquad\text{(even under parity)}.}
\]
If $A$ is \emph{odd}, $A'=-A$,
and the same manipulation gives
$A\PP=-\PP A$. Thus
\[
\boxed{\{\PP,A\}\equiv\PP A+A\PP=0
\qquad\text{(odd under parity)}.}
\]
Consequently a true scalar and
a pseudovector commute with parity;
a pseudoscalar and a true vector
anticommute with parity.
\phys{The scalar/vector distinction
concerns rotations; the true/pseudo
distinction concerns inversion.}

\problem{6.10}{Position and momentum are parity-odd}
Apply the triple product $\PP \x\PP$
to an arbitrary $f$:
\begin{align}
(\PP\x\PP f)(x)
&=\PP[x f(-x)]\\
&=-x f(x)=-(\x f)(x).
\end{align}
Therefore $\boxed{\PP\x\PP=-\x}$.
For momentum, using $\p=-\ii\hb\partial_x$,
\begin{align}
(\PP\p\PP f)(x)
&=\PP\left[-\ii\hb\frac{\dd}{\dd x}
f(-x)\right]\\
&=\PP[\ii\hb f'(-x)]
=\ii\hb f'(x)=-(\p f)(x).
\end{align}
Thus $\boxed{\PP\p\PP=-\p}$.
In three dimensions, apply the argument
to each coordinate:
\begin{equation}
\boxed{\PP^\dagger\mathbf r\PP=-\mathbf r,\qquad
\PP^\dagger\mathbf p\PP=-\mathbf p.}
\end{equation}
Their cross product
$\mathbf L=\mathbf r\times\mathbf p$
is consequently parity \emph{even}.

\problem{6.11}{Selection rule for a pseudovector or true scalar}
Let the two states have definite parities
$\PP\ket a=\eta_a\ket a$ and
$\PP\ket b=\eta_b\ket b$, where
$\eta_a,\eta_b=\pm1$.
For angular momentum, $\PP\mathbf L\PP=\mathbf L$.
Insert the identity as parity operators
on both sides:
\[
\bra b\mathbf L\ket a=
\bra b\PP^\dagger\mathbf L\PP\ket a
=\eta_b\eta_a\bra b\mathbf L\ket a.
\]
If $\eta_b\eta_a=-1$, the only possible
matrix element is zero:
\begin{equation}
\boxed{\bra{n'\ell'm'}\mathbf L\ket{n\ell m}=0
\quad\text{if }\ell+\ell'\text{ is odd}.}
\end{equation}
The same \emph{parity} rule holds for
any pseudovector or true scalar.
For $\mathbf L$ an even stronger
\emph{rotational} rule applies:
$\ell'=\ell$, because
$[\mathbf L,L^2]=0$ and $\mathbf L$
acts inside a fixed $\ell$ irreducible
multiplet.
\caveat{Parity gives a necessary condition,
not a guarantee of a nonzero matrix element.}

\problem{6.12}{Parity acts trivially on a pure spin-half spinor}
Let $\PP_s$ denote the $2\times2$ matrix
that acts only on the spin factor.
Spin is a pseudovector, so
$[\PP_s,S_i]=0$ for $i=x,y,z$.
Write
$\PP_s=\begin{pmatrix}a&b\\c&d\end{pmatrix}$.
From $[\PP_s,S_z]=0$ we obtain
$b=c=0$, so $\PP_s=\mathrm{diag}(a,d)$.
Then $[\PP_s,S_x]=0$ forces $a=d$:
$\PP_s=aI_2$.
The parity identity $\PP_s^2=I_2$
and Hermiticity give $a=\pm1$.
Choose the conventional intrinsic
spin parity $+1$:
\begin{equation}
\boxed{\PP_s=I_2.}
\end{equation}
The full electron state can still have
odd spatial parity $(-1)^\ell$.
This convention does not claim that an
electron and its antiparticle have the
same intrinsic parity.

\problem{6.13}{Electric dipole of the $1s$ state and an $n=2$ superposition}
\textbf{(a)} For an electron
$\mathbf p_e=-e\mathbf r$,
a parity-odd vector.
The hydrogen $1s$ state has even parity,
so $\av{\mathbf p_e}=-
\av{\PP^\dagger\mathbf p_e\PP}
=-\av{\mathbf p_e}$.
Therefore
$\boxed{\av{\mathbf p_e}_{1s}=0}$,
with no integration.
\textbf{(b)} The $n=2$ shell is
degenerate between $2s$ and $2p$.
Take the \emph{normalized} state
\begin{equation}
\boxed{\ket\psi=
(\ket{200}+\ket{210})/\sqrt2.}
\end{equation}
Its $x$ and $y$ dipoles vanish by
azimuthal integration. For $z=r\cos\theta$,
the two diagonal terms vanish by parity,
and only the cross term survives:
\[
\av z=\frac12
(\bra{200}z\ket{210}
+\bra{210}z\ket{200}).
\]
Use the real hydrogen functions
\[
R_{20}=\frac{(2-r/a_0)\ee^{-r/(2a_0)}}{
2\sqrt2\,a_0^{3/2}},\quad
R_{21}=\frac{(r/a_0)\ee^{-r/(2a_0)}}{
2\sqrt6\,a_0^{3/2}},
\]
and $Y_0^0=1/\sqrt{4\pi}$,
$Y_1^0=\sqrt{3/(4\pi)}\cos\theta$.
The angular integral is
$\int (Y_0^0)^*\cos\theta Y_1^0\dd\Omega
=1/\sqrt3$. The radial one is
\begin{align}
\int_0^\infty r^3R_{20}R_{21}\dd r
&=\frac{a_0}{8\sqrt3}
\int_0^\infty t^4(2-t)\ee^{-t}\dd t\\
&=\frac{a_0}{8\sqrt3}(2\cdot4!-5!)
=-3\sqrt3\,a_0.
\end{align}
Hence $\bra{200}z\ket{210}=-3a_0$,
and
\begin{equation}
\boxed{\av{\mathbf r}=-3a_0\hat{\mathbf z},
\qquad\av{\mathbf p_e}=+3ea_0\hat{\mathbf z}.}
\end{equation}
\phys{A degenerate superposition of
opposite-parity orbitals can carry a
static dipole moment even though each
pure orbital has zero dipole.}

\problem{6.14}{Finite rotations from infinitesimal planar matrices}
Put $\delta=\varphi/N$ and
\[
M_N=\begin{pmatrix}1&-\delta\\\delta&1\end{pmatrix}
=I+\delta J,\qquad
J=\begin{pmatrix}0&-1\\1&0\end{pmatrix}.
\]
\textbf{(a)} Since $J$ has eigenvectors
$v_-=(1,\ii)^T/\sqrt2$ and
$v_+=(1,-\ii)^T/\sqrt2$
with eigenvalues $-\ii$ and $+\ii$,
the unitary $S=(v_-\ v_+)$ gives
\[
\boxed{S^\dagger M_NS=
\operatorname{diag}(1-\ii\delta,1+\ii\delta).}
\]
\textbf{(b)} For each eigenvalue,
$\lim_{N\to\infty}
(1\mp\ii\varphi/N)^N=\ee^{\mp\ii\varphi}$.
Thus
\[
\lim_{N\to\infty}M_N^N=
S\,\operatorname{diag}(\ee^{-\ii\varphi},
\ee^{+\ii\varphi})S^\dagger.
\]
\textbf{(c)} Evaluate the product:
\begin{equation}
\boxed{\lim_{N\to\infty}M_N^N=
\ee^{\varphi J}=
\begin{pmatrix}
\cos\varphi&-\sin\varphi\\
\sin\varphi&\cos\varphi
\end{pmatrix}.}
\end{equation}
This exactly connects the
infinitesimal transformation
of Equation 6.30 with its finite
rotation in Equation 6.31.

\problem{6.15}{Rotational invariance of a scalar operator}
A scalar operator satisfies
$[L_i,f]=0$ for each Cartesian $i$.
Hence $[\mathbf n\cdot\mathbf L,f]=0$
for every unit $\mathbf n$.
Because
$\RR_{\mathbf n}(\varphi)
=\exp(-\ii\varphi\,\mathbf n\cdot\mathbf L/\hb)$,
every power in its exponential commutes
with $f$ and so does the complete rotation:
$\RR f=f\RR$.
Consequently
\begin{equation}
\boxed{\RR^\dagger f\RR
=\RR^\dagger\RR f=f.}
\end{equation}
The result holds for both true scalars
and pseudoscalars under \emph{proper}
rotations. Parity is a separate question.

\problem{6.16}{Infinitesimal rotation about the $y$ axis}
For a vector operator $\mathbf V$,
$[L_i,V_j]=\ii\hb\epsilon_{ijk}V_k$.
Under a small rotation of angle
$\delta$ about $y$,
\begin{align}
V_j'&=\RR_y^\dagger(\delta)V_j\RR_y(\delta)\\
&=V_j+\frac{\ii\delta}{\hb}[L_y,V_j]
+O(\delta^2)\\
&=V_j-\delta\epsilon_{yjk}V_k+O(\delta^2).
\end{align}
Using $\epsilon_{yxz}=-1$,
$\epsilon_{yzx}=+1$:
\[
V_x'=V_x+\delta V_z,\quad
V_y'=V_y,\quad
V_z'=V_z-\delta V_x.
\]
Thus
\begin{equation}
\boxed{\begin{pmatrix}V_x'\\V_y'\\V_z'\end{pmatrix}
=\begin{pmatrix}
1&0&\delta\\
0&1&0\\
-\delta&0&1
\end{pmatrix}
\begin{pmatrix}V_x\\V_y\\V_z\end{pmatrix}
+O(\delta^2).}
\end{equation}
This is the linear approximation to
the conventional $y$-axis rotation
matrix.

\problem{6.17}{Infinitesimal rotation of $\ket{n\ell m}$}
Expand the rotation operator to first order:
\[
\RR_{\mathbf n}(\delta)=
I-\frac{\ii\delta}{\hb}
(n_xL_x+n_yL_y+n_zL_z)+O(\delta^2).
\]
Use $L_x=(L_++L_-)/2$,
$L_y=(L_+-L_-)/(2\ii)$, giving
\[
\mathbf n\cdot\mathbf L=
n_zL_z+\frac{n_x-\ii n_y}{2}L_+
+\frac{n_x+\ii n_y}{2}L_-.
\]
On the angular-momentum basis,
\[
L_\pm\ket{n\ell m}=
\hb\sqrt{(\ell\mp m)(\ell\pm m+1)}
\ket{n\ell,m\pm1},
\quad L_z\ket{n\ell m}=m\hb\ket{n\ell m}.
\]
Therefore
\begin{align}
\RR_{\mathbf n}(\delta)\ket{n\ell m}
&=\sum_{m'=-\ell}^{\ell}
D_{m'm}^{(\ell)}\ket{n\ell m'}+O(\delta^2),\\
\boxed{D_{m'm}^{(\ell)}}&=
\delta_{m'm}-\ii\delta\bigg[
n_zm\,\delta_{m'm}\\
&\quad+\frac{n_x-\ii n_y}{2}
\sqrt{(\ell-m)(\ell+m+1)}\,
\delta_{m',m+1}\\
&\quad+\frac{n_x+\ii n_y}{2}
\sqrt{(\ell+m)(\ell-m+1)}\,
\delta_{m',m-1}
\bigg].
\end{align}
The principal $n$ and total orbital
$\ell$ labels remain fixed, while
the projection $m$ mixes with its
neighbors. As $D=I-\ii\delta A$
with $A=A^\dagger$, $D^\dagger D
=I+O(\delta^2)$, as required
by unitarity.

\section{Degeneracy and tensor selection rules: Problems 6.18--6.25}
\problem{6.18}{Why translation and parity generate free-particle degeneracy}
\textbf{(a)} For any test function $f$,
\[
(\PP\T(a)f)(x)=f(-x-a),\qquad
(\T(a)\PP f)(x)=f(a-x).
\]
These are not generally equal; instead
\begin{equation}
\boxed{\PP\T(a)=\T(-a)\PP.}
\end{equation}
\textbf{(b)} The momentum eigenfunction
$f_p(x)=\ee^{\ii px/\hb}/\sqrt{2\pi\hb}$
obeys
\[
\PP f_p(x)=f_p(-x)=f_{-p}(x).
\]
The free Hamiltonian $H=\p^2/(2m)$
commutes with parity. Thus if
$Hf_p=E_pf_p$,
$H(\PP f_p)=\PP Hf_p=E_p f_{-p}$.
Hence
\begin{equation}
\boxed{E_{-p}=E_p=p^2/(2m),\qquad p\ne0.}
\end{equation}
\textbf{(c)} Define the even and odd
continuum-normalized functions
$c_p(x)=\cos(px/\hb)/\sqrt{\pi\hb}$,
$s_p(x)=\sin(px/\hb)/\sqrt{\pi\hb}$.
Angle-addition formulae give
\[
\begin{pmatrix}\T(a)c_p\\\T(a)s_p\end{pmatrix}
=
\begin{pmatrix}
\cos(pa/\hb)&\sin(pa/\hb)\\
-\sin(pa/\hb)&\cos(pa/\hb)
\end{pmatrix}
\begin{pmatrix}c_p\\s_p\end{pmatrix}.
\]
They are mixed by translations, so
they must belong to the same energy
eigenspace if translations commute
with $H$.
\phys{The momentum and parity bases
diagonalize different commuting
pairs; their ability to mix at fixed
energy reveals the degeneracy.}

\problem{6.19}{Commutators of a vector ladder operator}
Let $V_\pm=V_x\pm\ii V_y$,
and assume vector commutators
$[L_i,V_j]=\ii\hb\epsilon_{ijk}V_k$.
\textbf{(a)} Direct substitution yields
\[
[L_z,V_x]=\ii\hb V_y,\qquad
[L_z,V_y]=-\ii\hb V_x,
\]
so
\begin{equation}
\boxed{[L_z,V_\pm]=\pm\hb V_\pm.}
\end{equation}
Use $L^2=L_-L_++L_z^2+\hb L_z$
and
\[
[L_+,V_+]=0,\quad
[L_-,V_+]=-2\hb V_z,\quad
[L_z,V_+]=\hb V_+.
\]
Expanding $[AB,C]=A[B,C]+[A,C]B$
and commuting $L_z$ through the
vector components gives
\begin{equation}
\boxed{[L^2,V_\pm]=
2\hb^2V_\pm
\pm2\hb V_\pm L_z
\mp2\hb V_zL_\pm.}
\end{equation}
\textbf{(b)} Suppose
$L^2\ket{\ell,\ell}
=\hb^2\ell(\ell+1)\ket{\ell,\ell}$.
Since $L_+\ket{\ell,\ell}=0$,
the previous formula gives
\[
L^2 V_+\ket{\ell,\ell}
=\hb^2[\ell(\ell+1)+2+2\ell]
V_+\ket{\ell,\ell}.
\]
The factor is $(\ell+1)(\ell+2)$.
Also
$L_zV_+\ket{\ell,\ell}
=(\ell+1)\hb V_+\ket{\ell,\ell}$.
Therefore either
$V_+\ket{\ell,\ell}=0$, or it has
\begin{equation}
\boxed{\ell'=\ell+1,\qquad
m'=\ell+1.}
\end{equation}
This is the highest-weight special
case of the rank-one tensor
selection rules.

\problem{6.20}{Scalar selection rule using the lowering operator}
A scalar $f$ has $[L_-,f]=0$.
Fix $\ell'=\ell$ and consider
$\ket{n\ell m}$ and $\ket{n'\ell,m-1}$.
The matrix element of the commutator is
\[
0=\bra{n'\ell,m-1}[L_-,f]
\ket{n\ell m}.
\]
For a scalar, $[L_z,f]=0$, so only
equal magnetic indices occur.
Using
$L_-\ket{\ell m}=
\hb B^\ell_m\ket{\ell,m-1}$,
$B^\ell_m=\sqrt{(\ell+m)(\ell-m+1)}$,
we obtain
\[
0=\hb B^\ell_m
\left(
\bra{n'\ell m}f\ket{n\ell m}
-\bra{n'\ell,m-1}f\ket{n\ell,m-1}
\right).
\]
For $m>-\ell$ the prefactor is
nonzero, giving
\begin{equation}
\boxed{\bra{n'\ell m}f\ket{n\ell m}
=\bra{n'\ell,m-1}f\ket{n\ell,m-1}.}
\end{equation}
Thus the diagonal-in-$m$ matrix element
is independent of $m$, the same
selection relation obtained from
$[L_+,f]=0$.

\problem{6.21}{A scalar reduced matrix element for a hydrogen superposition}
Take the normalized
$\ket\psi=(\ket{211}+\ket{21,-1})/\sqrt2$.
Rotational scalar selection rules
for the magnitude $r$ state
\[
\bra{n'\ell'm'}r\ket{n\ell m}
=\delta_{\ell'\ell}\delta_{m'm}
\mathcal R_{n'\ell;n\ell},
\]
where $\mathcal R_{n'\ell;n\ell}$
is the single reduced radial
matrix element for fixed $\ell$.
The cross terms have $m'\ne m$ and
vanish. The diagonal ones are equal,
so
\[
\av r=\frac12(\mathcal R_{21;21}
+\mathcal R_{21;21})=\mathcal R_{21;21}.
\]
For hydrogen's $2p$ radial function
$R_{21}(r)=
r\ee^{-r/(2a_0)}/(2\sqrt6\,a_0^{5/2})$,
\begin{align}
\mathcal R_{21;21}
&=\int_0^\infty
r^3|R_{21}(r)|^2\dd r\\
&=\frac1{24a_0^5}
\int_0^\infty r^5\ee^{-r/a_0}\dd r
=\frac{5!a_0^6}{24a_0^5}.
\end{align}
Consequently
\begin{equation}
\boxed{\av r=5a_0.}
\end{equation}
\phys{Rotational symmetry reduces two
apparently separate diagonal
expectation values to one radial integral.}

\problem{6.22}{Vector-operator commutators and magnetic selection rule}
\textbf{(a)} Start from
$[L_i,V_j]=\ii\hb\epsilon_{ijk}V_k$.
Using $L_\pm=L_x\pm\ii L_y$
and $V_\pm=V_x\pm\ii V_y$,
expand the commutators componentwise.
The results are
\begin{align}
\boxed{[L_z,V_z]}&=0,&
\boxed{[L_z,V_\pm]}&=\pm\hb V_\pm,\\
\boxed{[L_\pm,V_\pm]}&=0,&
\boxed{[L_\pm,V_z]}&=\mp\hb V_\pm,\\
\boxed{[L_\pm,V_\mp]}&=\pm2\hb V_z.
\end{align}
For example
$[L_+,V_z]=[L_x,V_z]
+\ii[L_y,V_z]=-\ii\hb V_y-\hb V_x
=-\hb V_+$; the remaining identities
follow analogously.
\textbf{(b)} Insert eigenstates
$\ket a=\ket{n\ell m}$ and
$\ket b=\ket{n'\ell'm'}$
into $[L_z,V_z]=0$:
\[
0=\bra b[L_z,V_z]\ket a
=\hb(m'-m)\bra b V_z\ket a.
\]
Therefore
\begin{equation}
\boxed{\bra{n'\ell'm'}V_z\ket{n\ell m}
=0\quad\text{unless }m'=m.}
\end{equation}
Likewise $V_+$ needs $m'=m+1$,
and $V_-$ needs $m'=m-1$.
These are the magnetic quantum-number
selection rules in Equations 6.56--6.58.

\problem{6.23}{Clebsch--Gordan coefficient recursions}
\textbf{(a)} In the product space
$\ket{j_1j_2m_1m_2}$,
\[
\ket{JM}=\sum_{m_1+m_2=M}
C_{m_1m_2;M}^{j_1j_2;J}
\ket{j_1j_2m_1m_2}.
\]
Take the inner product with a fixed
uncoupled basis vector.
Orthonormality leaves exactly one term:
\begin{equation}
\boxed{C_{m_1m_2;M}^{j_1j_2;J}
=\braket{j_1j_2m_1m_2}{JM}.}
\end{equation}
\textbf{(b)} Define ladder coefficients
\[
A_m^j=\sqrt{(j-m)(j+m+1)},\qquad
B_m^j=\sqrt{(j+m)(j-m+1)}.
\]
Apply $J_+=J_{1+}+J_{2+}$ to the
coupled decomposition of $\ket{JM}$.
The left-hand side is
$\hb A_M^J\ket{J,M+1}$.
Match the coefficient of the
uncoupled $\ket{m_1,m_2}$:
\begin{equation}
\boxed{
A_M^J C^{J,M+1}_{m_1m_2}=
B_{m_1}^{j_1}C^{J,M}_{m_1-1,m_2}
+B_{m_2}^{j_2}C^{J,M}_{m_1,m_2-1}.}
\end{equation}
Similarly, $J_-$ yields
\begin{equation}
\boxed{
B_M^J C^{J,M-1}_{m_1m_2}=
A_{m_1}^{j_1}C^{J,M}_{m_1+1,m_2}
+A_{m_2}^{j_2}C^{J,M}_{m_1,m_2+1}.}
\end{equation}
In these formulas the fixed superscripts
$j_1,j_2,J$ on the $C$ coefficients
are suppressed for clarity.
Out-of-range magnetic indices have
zero coefficient. This establishes
both recursions directly from angular
momentum addition.

\problem{6.24}{Six vector matrix-element recursions and Wigner--Eckart}
Define
\[
M^{q}_{m'm}=\bra{n'\ell'm'}V_q\ket{n\ell m},
\quad q=+,-,z,
\]
and $A_m^\ell,B_m^\ell$
as in Problem 6.23.
\textbf{(a)} Sandwiching the six
commutators $[L_\pm,V_\pm]=0$,
$[L_\pm,V_z]=\mp\hb V_\pm$,
$[L_\pm,V_\mp]=\pm2\hb V_z$
between the angular eigenstates
produces the six identities
\begin{align}
B_{m'}^{\ell'}M^+_{m'-1,m}
-A_m^\ell M^+_{m',m+1}&=0,\\
A_{m'}^{\ell'}M^-_{m'+1,m}
-B_m^\ell M^-_{m',m-1}&=0,\\
B_{m'}^{\ell'}M^z_{m'-1,m}
-A_m^\ell M^z_{m',m+1}&=-M^+_{m'm},\\
A_{m'}^{\ell'}M^z_{m'+1,m}
-B_m^\ell M^z_{m',m-1}&=+M^-_{m'm},\\
B_{m'}^{\ell'}M^-_{m'-1,m}
-A_m^\ell M^-_{m',m+1}&=+2M^z_{m'm},\\
A_{m'}^{\ell'}M^+_{m'+1,m}
-B_m^\ell M^+_{m',m-1}&=-2M^z_{m'm}.
\end{align}
The terms on the left are obtained by
moving the ladder operator to the bra
or acting on the ket; the sign on the
right comes from its commutator.
\textbf{(b)} Adopt the spherical tensor
convention $V_{\pm1}=\mp V_\pm/\sqrt2$,
$V_0=V_z$. The Wigner--Eckart
ansatz (with phase convention of the text) is
\begin{equation}
\boxed{\begin{aligned}
M^+_{m'm}&=-\sqrt2\,
C^{\ell',m'}_{\ell,m;1,1}\,\mathcal R,\\
M^-_{m'm}&=+\sqrt2\,
C^{\ell',m'}_{\ell,m;1,-1}\,\mathcal R,\\
M^z_{m'm}&=
C^{\ell',m'}_{\ell,m;1,0}\,\mathcal R,
\end{aligned}}
\end{equation}
where $\mathcal R=\bra{n'\ell'}V\ket{n\ell}$
is a reduced matrix element independent of
$m,m'$ and tensor component.
Insert these three expressions into
the six commutator identities.
In each equation the common factor
$\mathcal R$ cancels, while the remaining
Clebsch--Gordan terms are precisely
the raising or lowering recursions
derived in Problem 6.23.
Hence the Wigner--Eckart form
solves all six constraints.
\phys{Rotations fix every dependence
on $m,m'$; radial dynamics is carried
only by the reduced matrix element.}

\problem{6.25}{The dipole of a $2p$ hydrogen superposition}
Take
\[
\ket\psi=(\ket{211}+\ket{210})/\sqrt2,\qquad
\mathbf p_e=-e\mathbf r.
\]
Each orbital has $\ell=1$, so
$\PP\ket\psi=-\ket\psi$;
the entire superposition has definite
odd parity. Because
$\PP^\dagger\mathbf p_e\PP=-\mathbf p_e$,
\[
\bra\psi\mathbf p_e\ket\psi
=\bra\psi\PP^\dagger\mathbf p_e\PP\ket\psi
=-\bra\psi\mathbf p_e\ket\psi.
\]
Therefore
\begin{equation}
\boxed{\av{p_{e,x}}=\av{p_{e,y}}=
\av{p_{e,z}}=0.}
\end{equation}
In Wigner--Eckart notation every
component is proportional to
the same reduced dipole matrix element
$\bra{2,1}\mathbf p_e\ket{2,1}$,
which is \emph{zero by Laporte's rule}.
Thus the vanishing is independent
of the Clebsch--Gordan numerical values;
no radial integration is required.

\section{Time translations, dynamics and propagators: Problems 6.26--6.31}
\problem{6.26}{Heisenberg momentum of the oscillator}
For the harmonic oscillator
$H=\p^2/(2m)+m\omega^2\x^2/2$,
Example 6.7 finds
\[
x_H(t)=x\cos\omega t+
\frac{p}{m\omega}\sin\omega t.
\]
The Heisenberg equations are
$\dot x_H=p_H/m$ and
$\dot p_H=-m\omega^2x_H$.
Differentiate the displayed
$x_H(t)$ and multiply by $m$:
\begin{equation}
\boxed{p_H(t)=p\cos\omega t
-m\omega x\sin\omega t.}
\end{equation}
At $t=0$ it reduces to $p$; its
time derivative is indeed
\[
\dot p_H=-\omega p\sin\omega t
-m\omega^2x\cos\omega t
=-m\omega^2x_H(t).
\]
These are exactly the \emph{operator}
forms of classical oscillator
trajectories, with $x,p$
replaced by their time-zero
noncommuting operators.

\problem{6.27}{Free-particle Heisenberg position and momentum}
For $H=p^2/(2m)$,
$[H,p]=0$ and
$[x,H]=\ii\hb p/m$.
Since the commutator $[x,H]$
commutes with $H$, the identity
\[
[x,H^n]=
\sum_{j=0}^{n-1}H^j[x,H]H^{n-1-j}
=n\frac{\ii\hb p}{m}H^{n-1}
\]
holds for every positive integer $n$.
Expand $U(t)=\sum_{n\ge0}(-\ii t/\hb)^nH^n/n!$:
\[
[x,U(t)]=
\sum_{n=1}^{\infty}
\frac{(-\ii t/\hb)^n}{n!}
[x,H^n]=\frac{t}{m}p\,U(t).
\]
Using $U^\dagger U=I$,
\begin{align}
x_H(t)&=U^\dagger xU
=x+U^\dagger[x,U]
=x+\frac{t}{m}p,\\
p_H(t)&=U^\dagger pU=p.
\end{align}
Thus
\begin{equation}
\boxed{x_H(t)=x_H(0)+
\frac{t}{m}p_H(0),\qquad
p_H(t)=p_H(0).}
\end{equation}
The mean values follow classical
uniform-velocity motion, while
the operators themselves still obey
$[x_H(t),p_H(t)]=\ii\hb$.

\problem{6.28}{Infinitesimal time-evolution operator}
Taylor-expand the state in
an infinitesimal interval $\delta$:
\[
\Psi(x,t_0+\delta)
=\Psi(x,t_0)+\delta\,\partial_t\Psi(x,t_0)
+O(\delta^2).
\]
The Schrödinger equation gives
$\partial_t\Psi=-(\ii/\hb)H(t_0)\Psi$.
Substitute:
\[
\Psi(x,t_0+\delta)=
\left[I-\frac{\ii\delta}{\hb}H(t_0)\right]
\Psi(x,t_0)+O(\delta^2).
\]
Therefore
\begin{equation}
\boxed{U(t_0+\delta,t_0)
=I-\frac{\ii\delta}{\hb}H(t_0)+O(\delta^2).}
\end{equation}
When $H$ depends on time, this
first-order expression is valid
without requiring $[H(t_1),H(t_2)]=0$.
A finite interval generally
requires time ordering.

\problem{6.29}{Heisenberg equations directly from the evolution operator}
Take a time-independent Hermitian
Hamiltonian and observable $Q$.
With $U(t)=\ee^{-\ii Ht/\hb}$,
$Q_H(t)=U^\dagger(t)QU(t)$.
Differentiate both factors:
\begin{align}
\frac{\dd Q_H}{\dd t}
&=\frac{\ii}{\hb}U^\dagger HQU
-\frac{\ii}{\hb}U^\dagger QHU\\
&=\frac{\ii}{\hb}[H,Q_H].
\end{align}
Equivalently
\begin{equation}
\boxed{\ii\hb\frac{\dd Q_H}{\dd t}
=[Q_H,H].}
\end{equation}
For a one-dimensional
$H=p^2/(2m)+V(x)$, use
$[x,p^2]=2\ii\hb p$ and
$[p,V(x)]=-\ii\hb V'(x)$.
This yields
\begin{equation}
\boxed{\frac{\dd x_H}{\dd t}=
\frac{p_H}{m},\qquad
\frac{\dd p_H}{\dd t}=
-\frac{\partial V(x_H)}{\partial x_H}.}
\end{equation}
The operator derivative of $V$
is interpreted through the
spectral functional calculus.
This is the Heisenberg analogue
of Hamilton's equations.

\problem{6.30}{Propagators for oscillator and free Gaussian packets}
\textbf{(a) Spectral kernel.}
Let $H\psi_n=E_n\psi_n$ and
$\int\psi_n^*(x')\Psi(x',0)\dd x'=c_n$.
The spectral resolution of $U$ is
$\sum_n\ket n\ee^{-\ii E_nt/\hb}\bra n$,
so
\begin{align}
\Psi(x,t)&=\sum_n\psi_n(x)
\ee^{-\ii E_nt/\hb}c_n\\
&=\int K(x,x';t)\Psi(x',0)\dd x',\\
\boxed{K(x,x';t)}&=
\sum_n\psi_n(x)\ee^{-\ii E_nt/\hb}
\psi_n^*(x').
\end{align}
At $t=0$, completeness gives
$K(x,x';0)=\delta(x-x')$.
\textbf{(b) Harmonic oscillator.}
Let $\theta=\omega t$,
$\xi=\sqrt{m\omega/\hb}\,x$,
$\eta=\sqrt{m\omega/\hb}\,x'$.
Insert the normalized Hermite
eigenfunctions
$\psi_n(x)\propto
\ee^{-\xi^2/2}H_n(\xi)/\sqrt{2^nn!}$
and $E_n=\hb\omega(n+\frac12)$
into the spectral sum.
The generating identity given in
the exercise (Mehler's formula),
with $z=\ee^{-\ii\theta}$,
sums it to
\begin{equation}
\boxed{
K_{\rm HO}(x,x';t)=
\sqrt{\frac{m\omega}{2\pi\ii\hb\sin\omega t}}
\exp\!\left\{
\frac{\ii m\omega}{2\hb\sin\omega t}
\big[(x^2+x'^2)\cos\omega t-2xx'\big]
\right\}.}
\end{equation}
This is valid away from focal
times $\sin\omega t=0$, with the
square-root phase chosen by
continuous evolution from $t=0$;
at those times the kernel is a
distribution with a Maslov phase.
\textbf{(c) Evolve an initially displaced Gaussian.}
Let
\[
\Psi(x',0)=\left(\frac{2a}{\pi}\right)^{1/4}
\ee^{-a(x'-x_0)^2}.
\]
Write
$B=m\omega/(2\hb\sin\theta)$,
$A=a-\ii B\cos\theta$, and
$D=2ax_0-2\ii Bx$.
Completing the square in the
integral $\int K_{\rm HO}\Psi_0\dd x'$
uses
\[
\int_{-\infty}^{\infty}
\ee^{-Ax'^2+Dx'}\dd x'
=\sqrt{\frac{\pi}{A}}\,
\ee^{D^2/(4A)},\qquad
\Re A=a>0.
\]
Consequently
\begin{equation}
\boxed{\begin{aligned}
\Psi(x,t)&=
\left(\frac{2a}{\pi}\right)^{1/4}
\sqrt{\frac{m\omega}{2\ii\hb\sin\theta\, A}}\\
&\quad\times\exp\left[
\ii Bx^2\cos\theta-ax_0^2+
\frac{D^2}{4A}\right].
\end{aligned}}
\end{equation}
This compact expression is exact and
includes the phase.
For $a=m\omega/(2\hb)$, its
probability density is a
minimum-width Gaussian centered at
$\boxed{\av{x(t)}=x_0\cos\omega t}$,
the textbook's special coherent-state case.
\textbf{(d) Free-particle kernel.}
Replace the discrete sum by an
integral of continuum momentum
eigenfunctions:
\[
K_0(x,x';t)=\int_{\mathbb R}
\frac{\dd p}{2\pi\hb}
\exp\left[\frac{\ii p(x-x')}{\hb}
-\frac{\ii p^2t}{2m\hb}\right].
\]
Complete the square in $p$
using the Gaussian integral with
the standard $t\to t-\ii0^+$
prescription:
\begin{equation}
\boxed{K_0(x,x';t)=
\sqrt{\frac{m}{2\pi\ii\hb t}}
\exp\left[\frac{\ii m(x-x')^2}{2\hb t}\right].}
\end{equation}
\textbf{(e) Free spreading Gaussian.}
With $\Psi(x,0)=(2a/\pi)^{1/4}\ee^{-ax^2}$,
another Gaussian integral gives
\begin{equation}
\boxed{\Psi(x,t)=
\left(\frac{2a}{\pi}\right)^{1/4}
\frac{1}{\sqrt{1+2\ii\hb at/m}}
\exp\left[-\frac{ax^2}{
1+2\ii\hb at/m}\right].}
\end{equation}
Set $\tau=2\hb at/m$ and take the
absolute square:
\[
|\Psi(x,t)|^2=
\sqrt{\frac{2a}{\pi}}\,
\frac{1}{\sqrt{1+\tau^2}}\,
\exp[-2ax^2/(1+\tau^2)].
\]
Hence
$\boxed{(\Delta x_t)^2=
(1+\tau^2)/(4a)}$,
which verifies normalized free
wavepacket spreading.
\caveat{Although the textbook informally
calls $|K|^2$ a position-transition
probability, a continuum position
eigenstate is not normalizable.
The actual probability density for a
normalized initial state is $|\Psi(x,t)|^2$.}

\problem{6.31}{Translation of a cusp state via spectral resolution}
The normalized cusp function is
$\psi(x)=\sqrt\lambda\,\ee^{-\lambda|x|}$:
\[
\int_{-\infty}^{\infty}|\psi(x)|^2\dd x
=2\lambda\int_0^\infty
\ee^{-2\lambda x}\dd x=1.
\]
The momentum-space transform is
\begin{align}
\phi(p)&=\frac{\sqrt\lambda}{\sqrt{2\pi\hb}}
\int_{\mathbb R}
\ee^{-\lambda|x|}\ee^{-\ii px/\hb}\dd x\\
&=\frac{\sqrt\lambda}{\sqrt{2\pi\hb}}
\left(\frac1{\lambda+\ii p/\hb}
+\frac1{\lambda-\ii p/\hb}\right)\\
&=\boxed{\sqrt{\frac{2\lambda^3}{\pi\hb}}\,
\frac1{\lambda^2+(p/\hb)^2}}.
\end{align}
Now use the spectral representation of
the translation operator:
\[
(\T(a)\psi)(x)=
\frac1{\sqrt{2\pi\hb}}\int_{\mathbb R}
\ee^{\ii p(x-a)/\hb}\phi(p)\dd p.
\]
It is exactly the inverse Fourier
transform of $\phi$ at $x-a$:
\begin{equation}
\boxed{(\T(a)\psi)(x)=
\sqrt\lambda\,\ee^{-\lambda|x-a|}.}
\end{equation}
This proof is valid despite the
undefined ordinary first derivative
at the cusp $x=0$.
\phys{The spectral definition of a
unitary operator is more general
than termwise Taylor expansion
on smooth wavefunctions.}

\section{Spin rotations and hidden spatial symmetries: Problems 6.32--6.34}
\problem{6.32}{SU(2) rotations of a spin-half particle}
Let $\boldsymbol\sigma=(\sigma_x,\sigma_y,\sigma_z)$
be the Pauli matrices.
\textbf{(a)} Their multiplication rule is
\[
\sigma_i\sigma_j=\delta_{ij}I+
\ii\epsilon_{ijk}\sigma_k.
\]
For ordinary vectors $\mathbf a,\mathbf b$,
sum $a_ib_j\sigma_i\sigma_j$ to obtain
\begin{equation}
\boxed{(\mathbf a\cdot\boldsymbol\sigma)
(\mathbf b\cdot\boldsymbol\sigma)
=(\mathbf a\cdot\mathbf b)I+
\ii(\mathbf a\times\mathbf b)\cdot\boldsymbol\sigma.}
\end{equation}
\textbf{(b)} For a unit vector $\mathbf n$,
the previous formula gives
$(\mathbf n\cdot\boldsymbol\sigma)^2=I$.
Thus its even powers are $I$ and odd
powers are $\mathbf n\cdot\boldsymbol\sigma$.
Expand the exponential:
\begin{align}
R_{\mathbf n}(\varphi)
&=\exp[-\ii\varphi
\mathbf n\cdot\mathbf S/\hb]\\
&=\exp[-\ii(\varphi/2)
\mathbf n\cdot\boldsymbol\sigma]\\
&=\boxed{\cos(\varphi/2)I-
\ii\sin(\varphi/2)
\,\mathbf n\cdot\boldsymbol\sigma.}
\end{align}
\textbf{(c)} With
$\mathbf n=(\sin\theta\cos\phi,
\sin\theta\sin\phi,\cos\theta)$,
\[
\mathbf n\cdot\boldsymbol\sigma=
\begin{pmatrix}
\cos\theta&\sin\theta\,\ee^{-\ii\phi}\\
\sin\theta\,\ee^{\ii\phi}&-\cos\theta
\end{pmatrix}.
\]
Substitute this matrix into part (b)
to obtain the requested standard-basis
$2\times2$ rotation matrix.
\textbf{(d)} Let $N=\mathbf n\cdot\boldsymbol\sigma$,
which is Hermitian and satisfies $N^2=I$.
Then
\[
R^\dagger R=(cI+\ii sN)(cI-\ii sN)
=(c^2+s^2)I=I
\]
with $c=\cos(\varphi/2)$,
$s=\sin(\varphi/2)$.
Thus $\boxed{R^\dagger R=I}$.
\textbf{(e)} For a $z$ rotation,
$R_z=\operatorname{diag}(\ee^{-\ii\varphi/2},
\ee^{+\ii\varphi/2})$. Multiply explicitly:
\begin{align}
R_z^\dagger S_xR_z
&=\frac{\hb}{2}
\begin{pmatrix}0&\ee^{\ii\varphi}\\
\ee^{-\ii\varphi}&0\end{pmatrix}\\
&=\boxed{S_x\cos\varphi-S_y\sin\varphi}.
\end{align}
This agrees with the transformed
vector component in the textbook's
active-state convention.
\textbf{(f)} Around $x$ through $\pi$,
\[
\boxed{R_x(\pi)=-\ii\sigma_x=
\begin{pmatrix}0&-\ii\\-\ii&0\end{pmatrix}.}
\]
It sends $\ket{\uparrow_z}=(1,0)^T$
to $-\ii\ket{\downarrow_z}$, so the
spin is flipped up to an
unobservable overall phase.
\textbf{(g)} A complete $2\pi$ turn gives
\begin{equation}
\boxed{R_z(2\pi)=
\operatorname{diag}(-1,-1)=-I.}
\end{equation}
A spinor changes sign under $2\pi$
rotation; it returns to itself only
after $4\pi$. The sign is physically
observable through interference
against an unrotated reference path.

\problem{6.33}{Dihedral symmetry and a degenerate square-well pair}
Use the centered square $-L/2<x,y<L/2$,
and write
$f_n(x)=\sqrt{2/L}\sin[n\pi(x/L-1/2)]$,
$\psi_{ab}(x,y)=f_a(x)f_b(y)$.
\textbf{(a)} Under active counterclockwise
rotation by $\pi/2$, the wavefunction
is evaluated at the inverse-rotated
coordinates $(y,-x)$:
\[
(\RR_{\pi/2}\psi_{ab})(x,y)
=f_a(y)f_b(-x).
\]
An individual centered box mode
satisfies
\[
f_n(-x)=(-1)^{n+1}f_n(x).
\]
Therefore
\begin{equation}
\boxed{\RR_{\pi/2}\psi_{ab}
=(-1)^{b+1}\psi_{ba},\qquad
\RR_{\pi/2}\psi_{ba}
=(-1)^{a+1}\psi_{ab}.}
\end{equation}
\textbf{(b)} If $a$ and $b$ have
the same parity, then
$s=(-1)^{a+1}=(-1)^{b+1}$.
Define $\psi_\pm=(\psi_{ab}\pm\psi_{ba})/\sqrt2$.
Using part (a),
\[
\RR_{\pi/2}\psi_\pm=
s(\psi_{ba}\pm\psi_{ab})/\sqrt2,
\]
so
\begin{equation}
\boxed{\RR_{\pi/2}\psi_+=s\psi_+,\qquad
\RR_{\pi/2}\psi_-=-s\psi_-.}
\end{equation}
For two odd indices, $s=+1$;
for two even indices, $s=-1$.
\textbf{(c)} For $a=5$, $b=7$,
both one-dimensional modes are even
under $x\mapsto-x$.
Hence reflections in the
horizontal or vertical axis preserve
$\psi_-$. Reflection in the
diagonal $x=y$ interchanges the
two products and changes the sign:
\[
\psi_-(y,x)=-\psi_-(x,y).
\]
The other diagonal combines a
sign-preserving axial reflection
with this interchange, so it also
gives a minus sign.
A quarter turn gives a minus sign
by (b), and its powers produce
$\pm1$ factors. Therefore $\psi_-$
is an eigenstate of \emph{every}
symmetry in the square's dihedral
group $D_4$.
The script writes a contour plot of
$\psi_-=\frac1{\sqrt2}(\psi_{57}-\psi_{75})$,
whose nodal and sign pattern visibly
obeys these transformations.
\phys{Although $\psi_{57}$ and
$\psi_{75}$ have the same energy
$E\propto5^2+7^2$, their symmetric
and antisymmetric combinations
transform in inequivalent one-dimensional
$D_4$ representations. This discrete
symmetry alone does not \emph{force}
their equality of energy; the
separable square Hamiltonian has
additional structure.}

\problem{6.34}{Laplace--Runge--Lenz operator and hidden hydrogen degeneracy}
Put $k=e^2/(4\pi\epsilon_0)$.
The Hermitian quantum
Laplace--Runge--Lenz vector is
\[
\mathbf M=\frac{\mathbf p\times\mathbf L-
\mathbf L\times\mathbf p}{2m}
-\frac{k}{r}\,\mathbf r.
\]
The exercise supplies the algebra
$[H,M_i]=0$, $[L_i,M_j]=
\ii\hb\epsilon_{ijk}M_k$,
$[M_i,M_j]=-(2\ii\hb H/m)
\epsilon_{ijk}L_k$, together with
\begin{equation}
\mathbf M^2=k^2+\frac{2H}{m}
(\mathbf L^2+\hb^2),
\qquad\mathbf M\cdot\mathbf L=0.
\end{equation}
The latter are taken as given,
as specified by the exercise.
Write $\ket{n\ell\ell}$ for a
hydrogen state with maximal $m$.
Rotational covariance (Problem 6.19)
and $[H,M_+]=0$ imply
\[
M_+\ket{n\ell\ell}
=c_{n\ell}\ket{n,\ell+1,\ell+1}
\]
unless the state is annihilated.
\textbf{(a)} The identity
$\mathbf M\cdot\mathbf L=0$ can be
written
\[
M_zL_z+\frac12(M_+L_-+M_-L_+)=0.
\]
Apply it to the highest-weight
state, for which $L_+\ket{n\ell\ell}=0$.
Use $L_-\ket{\ell\ell}
=\hb\sqrt{2\ell}\ket{\ell,\ell-1}$.
The Wigner--Eckart ratio of relevant
rank-one tensor matrix elements gives
\[
M_+\ket{n\ell,\ell-1}
=c_{n\ell}\sqrt{\frac{\ell}{\ell+1}}
\ket{n,\ell+1,\ell},
\]
which follows from the ordinary
CG coefficients for $\ell\otimes1
\to\ell+1$.
For $\ell>0$ substitute to obtain
\begin{equation}
\boxed{M_z\ket{n\ell\ell}
=-\frac{c_{n\ell}}{\sqrt{2(\ell+1)}}
\ket{n,\ell+1,\ell}.}
\end{equation}
The $\ell=0$ case follows by
rotational covariance / CG factors,
since the scalar product relation
alone becomes trivial there.
\textbf{(b)} With
$M_\pm=M_x\pm\ii M_y$,
\[
M_-M_+=M_x^2+M_y^2+
\ii[M_x,M_y]
=M^2-M_z^2+\frac{2\hb H}{m}L_z.
\]
Act on $\ket{n\ell\ell}$ and insert
$E_n=-mk^2/(2\hb^2n^2)$,
$\mathbf L^2=\hb^2\ell(\ell+1)$ and
$L_z=\hb\ell$:
\begin{align}
M_-M_+\ket{n\ell\ell}
&=\left[k^2-\frac{k^2}{n^2}
(\ell(\ell+1)+1)
-\frac{k^2\ell}{n^2}-M_z^2\right]
\ket{n\ell\ell}\\
&=\boxed{\left[
k^2\left(1-\frac{(\ell+1)^2}{n^2}\right)
-M_z^2\right]\ket{n\ell\ell}.}
\end{align}
\textbf{(c)} Take the expectation
value in the normalized state.
Because $M_-=M_+^\dagger$,
the left side is
$\|M_+\ket{n\ell\ell}\|^2=|c_{n\ell}|^2$.
From part (a),
$\|M_z\ket{n\ell\ell}\|^2=
|c_{n\ell}|^2/[2(\ell+1)]$.
Therefore
\[
|c_{n\ell}|^2=
k^2\left(1-\frac{(\ell+1)^2}{n^2}\right)
-\frac{|c_{n\ell}|^2}{2(\ell+1)}.
\]
Solving for the coefficient:
\begin{equation}
\boxed{|c_{n\ell}|=
k\sqrt{\frac{2(\ell+1)}{2\ell+3}
\left[1-\frac{(\ell+1)^2}{n^2}\right]}.}
\end{equation}
An arbitrary phase may be chosen
for the states, so only the magnitude
is physically fixed.
This coefficient is nonzero for
$0\le\ell\le n-2$, but vanishes
for $\ell=n-1$, where no state
with $\ell=n$ exists in the same
hydrogen shell.
\phys{The LRL vector moves between
different $\ell$ multiplets at fixed
principal quantum number $n$.
Together with angular momentum,
it generates the hidden SO(4)
symmetry behind the $n^2$
orbital degeneracy of hydrogen
(with spin ignored).}

\section{Galilean boosts and time reversal: Problems 6.35--6.37}
\problem{6.35}{A Galilean boost as a unitary quantum transformation}
Define
\[
G(v,t)=\exp\left[
\frac{\ii}{\hb}(mv\x-vt\p)\right]
\]
with the textbook's active-state
boost convention.
\textbf{(a)} For infinitesimal $v=\delta$,
write $G=I+A+O(\delta^2)$,
$A=\ii\delta(mx-tp)/\hb$.
As $G^\dagger xG=x+[x,A]+O(\delta^2)$,
the canonical commutator gives
\[
[x,A]=\delta t,\qquad
[p,A]=m\delta.
\]
Therefore
\begin{equation}
\boxed{G^\dagger xG=x+\delta t,\qquad
G^\dagger pG=p+m\delta+O(\delta^2).}
\end{equation}
The wave packet's position and
momentum are shifted according to
a frame moving with relative speed $v$.
\textbf{(b)} Let
$A=\ii mvx/\hb$, $B=-\ii vt p/\hb$.
Their commutator is the scalar
$[A,B]=\ii mv^2t/\hb$, and all
nested commutators vanish.
The Baker--Campbell--Hausdorff
identity therefore yields
\begin{align}
G&=\ee^{A+B}
=\ee^{A}\ee^{B}\ee^{-[A,B]/2}\\
&=\boxed{
\exp\left[\frac{\ii}{\hb}
(mvx-\tfrac12mv^2t)\right]\T(vt)},\\
G&=\ee^{B}\ee^{A}\ee^{+[A,B]/2}\\
&=\boxed{\T(vt)
\exp\left[\frac{\ii}{\hb}
(mvx+\tfrac12mv^2t)\right]}.
\end{align}
\textbf{(c)} If $\Psi$ satisfies
$\ii\hb\partial_t\Psi=
[-\hb^2\partial_x^2/(2m)+V(x)]\Psi$,
then the boosted wave is
\begin{equation}
\boxed{\Psi'(x,t)=
\ee^{\ii(mvx-mv^2t/2)/\hb}\,
\Psi(x-vt,t).}
\end{equation}
Set $y=x-vt$ and denote the phase
by $F$. The time derivative is
\[
\ii\hb\partial_t\Psi'
=F[\tfrac12mv^2\Psi-\ii\hb v\partial_y\Psi+
\ii\hb\partial_t\Psi].
\]
Differentiate twice in $x$:
\[
-\frac{\hb^2}{2m}\partial_x^2\Psi'
=F[\tfrac12mv^2\Psi-\ii\hb v\partial_y\Psi
-\tfrac{\hb^2}{2m}\partial_y^2\Psi].
\]
Insert the original Schrödinger
equation; terms match term by term,
and the remaining potential is
$V(y)=V(x-vt)$:
\begin{equation}
\boxed{\ii\hb\partial_t\Psi'
=\left[-\frac{\hb^2}{2m}\partial_x^2
+V(x-vt)\right]\Psi'.}
\end{equation}
\textbf{(d)} This reproduces the
transformation of a wave packet
under a uniform Galilean boost
in Problem 2.50(a), including its
necessary space--time dependent
phase. Simply replacing
$x\mapsto x-vt$ without that phase
would not solve the moving-potential
Schrödinger equation.

\problem{6.36}{Spinless time reversal and the energy spectrum}
For a spinless particle in
position representation,
$\Theta=K$ (complex conjugation),
which is \emph{antiunitary}.
\textbf{(a)} Apply it to an arbitrary
test function $f$:
\[
(\Theta^{-1}x\Theta f)(x)=
[x f^*(x)]^*=xf(x).
\]
For momentum,
\[
(\Theta^{-1}p\Theta f)(x)
=[-\ii\hb (f^*)'(x)]^*
=+\ii\hb f'(x)=-pf(x).
\]
Thus
\begin{equation}
\boxed{\Theta^{-1}x\Theta=x,\qquad
\Theta^{-1}p\Theta=-p.}
\end{equation}
\textbf{(b)} Reversing motion should
undo unitary evolution:
\[
\Theta U(t)\Theta^{-1}=U(-t).
\]
Expand the operator
$U(\delta)=I-\ii H\delta/\hb+O(\delta^2)$.
Because $\Theta$ is antiunitary,
$\Theta\ii\Theta^{-1}=-\ii$.
Hence
\[
\Theta U(\delta)\Theta^{-1}
=I+\frac{\ii\delta}{\hb}
\Theta H\Theta^{-1}+O(\delta^2).
\]
Compare with
$U(-\delta)=I+\ii H\delta/\hb+O(\delta^2)$,
showing that time-reversal invariance is
equivalent to
\begin{equation}
\boxed{\Theta H\Theta^{-1}=H.}
\end{equation}
The relation is sometimes written
$[\Theta,H]=0$, but $\Theta$
is anti\emph{linear}; the equation
above is the precise statement.
\textbf{(c)} If $H\psi_n=E_n\psi_n$
with real $E_n$, then
$H\Theta\psi_n=\Theta H\psi_n
=E_n\Theta\psi_n$, so
$\psi_n^*$ has the same energy.
For a nondegenerate eigenvalue,
$\psi_n^*=c\psi_n$, $|c|=1$.
Choose a global phase
$\psi_n\mapsto\ee^{\ii\alpha}\psi_n$
to make $c=1$, proving that an
eigenfunction can be chosen real.
\textbf{(d)} Momentum eigenstates satisfy
\[
\boxed{\Theta f_p(x)=f_p^*(x)=f_{-p}(x)}.
\]
Both energies equal $p^2/(2m)$,
although their momenta are opposite.
For hydrogen
$\psi_{n\ell m}=R_{n\ell}(r)Y_\ell^m$,
using the real radial convention and
$Y_\ell^{m*}=(-1)^mY_\ell^{-m}$ gives
\begin{equation}
\boxed{\Theta\psi_{n\ell m}
=(-1)^m\psi_{n\ell,-m}.}
\end{equation}
Hydrogen energies depend only
on $n$ in the nonrelativistic
point-Coulomb model, so the
time-reversed state has the same
energy. For $m=0$ it is the same
state up to phase, while for
$m\ne0$ it is generally an
independent, degenerate state.
\phys{A fixed externally imposed
magnetic field breaks this symmetry
unless the magnetic field itself
is reversed with the system.}

\problem{6.37}{Time reversal of spin-half and Kramers degeneracy}
A spin-$1/2$ spinor is
$\chi=(a,b)^T$. With the
Condon--Shortley spin basis,
time reversal acts as
\begin{equation}
\boxed{\Theta\begin{pmatrix}a\\b\end{pmatrix}
=\begin{pmatrix}-b^*\\a^*\end{pmatrix}
=-\ii\sigma_yK\begin{pmatrix}a\\b\end{pmatrix}.}
\end{equation}
\textbf{(a)} Apply the operation a second time:
\[
\Theta^2\begin{pmatrix}a\\b\end{pmatrix}
=\Theta\begin{pmatrix}-b^*\\a^*\end{pmatrix}
=\begin{pmatrix}-a\\-b\end{pmatrix}.
\]
Therefore
\begin{equation}
\boxed{\Theta^2=-I\qquad(s=\tfrac12).}
\end{equation}
\textbf{(b)} Suppose a normalized energy
eigenstate $\ket\psi$ of a
time-reversal invariant Hamiltonian
were identical, up to phase,
to its time-reversed state:
$\Theta\ket\psi=c\ket\psi$.
As $\Theta$ is anti\emph{linear},
applying it again gives
\[
\Theta^2\ket\psi
=\Theta(c\ket\psi)
=c^*\Theta\ket\psi
=|c|^2\ket\psi.
\]
But part (a) demands
$\Theta^2\ket\psi=-\ket\psi$.
Thus $|c|^2=-1$, impossible.
Hence $\ket\psi$ and
$\Theta\ket\psi$ are distinct
eigenstates with the same energy:
\begin{equation}
\boxed{\text{Every level is at least doubly
degenerate if }\Theta^2=-I
\text{ and }[\Theta,H]=0.}
\end{equation}
The two states are in fact orthogonal.
Using antiunitarity
$\braket{\Theta a}{\Theta b}=
\braket b a$ with $a=\psi$, $b=\Theta\psi$,
\[
\braket{\Theta\psi}{\Theta^2\psi}
=\braket{\Theta\psi}{\psi}.
\]
The left side equals
$-\braket{\Theta\psi}{\psi}$,
so their overlap is zero.
This is Kramers degeneracy.
\caveat{The degeneracy theorem
applies to systems whose physical
time-reversal operator squares to
$-I$, such as an odd number of
spin-half fermions under the stated
assumptions. A generic spinless
particle has $\Theta^2=+I$ and
need not have a Kramers doublet.}

\section{Reproducible numerical verification}
The open-source script \texttt{chapter6/scripts/verify.py}
uses finite-difference spectra, analytic spin matrices,
discrete rotation operators and normalized Gaussian
integrals to check selected independent calculations.
\IfFileExists{figures/bloch_bands.pdf}{
\begin{figure}[htbp]\centering
\includegraphics[width=.82\linewidth]{figures/bloch_bands.pdf}
\caption{Lowest two bands of the periodic finite-depth
well used in Problem 6.6. Energies in $\hb^2/(2ma^2)$
and crystal momentum $qa$.}\end{figure}}{}
\IfFileExists{figures/square_symmetry.pdf}{
\begin{figure}[htbp]\centering
\includegraphics[width=.8\linewidth]{figures/square_symmetry.pdf}
\caption{Two-dimensional square-well antisymmetric
superposition $\psi_-=(\psi_{57}-\psi_{75})/\sqrt2$
(Problem 6.33).}\end{figure}}{}
\IfFileExists{figures/spin_rotation.pdf}{
\begin{figure}[htbp]\centering
\includegraphics[width=.8\linewidth]{figures/spin_rotation.pdf}
\caption{Spin-$1/2$ rotation probability check
over $0\le\varphi\le2\pi$ (Problem 6.32).}\end{figure}}{}
\clearpage
\section{References and limitations}
D. J. Griffiths and D. F. Schroeter,
\emph{Introduction to Quantum Mechanics}, 3rd edition,
Cambridge University Press (2018), Chapter 6,
\href{https://doi.org/10.1017/9781316995433}{DOI:
10.1017/9781316995433}.
Solutions are written for educational verification.
The author has not claimed ownership of the textbook
problems. The full statements, diagrams and specific
numerical conventions should be checked in the
textbook itself. Numerical studies specify
grid, units and convergence checks where applicable.
\end{document}